\documentclass[conference]{IEEEtran}
\IEEEoverridecommandlockouts

%%%%%%%%%%%%%%%%%%%%%%%%%%%%%%%%%%%%%%%%%%%%%%%%%%%%%%%%%%%%%%%%%%%%%%%%%%%%%%%%

% Math, plots and sketches
\usepackage{amsmath, mathtools, amsfonts, amssymb}
\usepackage{tkz-euclide}
\usepackage{pgfplots}
\usepackage{verbatim}
\pgfplotsset{compat=1.18}

% Images
\usepackage{graphicx}

% Allow line-breaking URLs/paths in text
\usepackage{url}
\def\UrlBreaks{\do\/\do-\do\_\do\#\do\0\do\1\do\2\do\3\do\4\do\5\do\6\do\7\do\8\do\9}

% Formatting
\usepackage{dirtytalk} % Proper quotation marks
\usepackage{float} % Placement [H]
\usepackage{algorithmic}
\usepackage{textcomp}

% Colors for highlighting and coloring tables
\usepackage{color, soul}
\usepackage[table]{xcolor}

% Custom labeled lists
\usepackage{enumerate}

% Clickable references
\usepackage{hyperref}
\hypersetup{
  colorlinks   = true, 
  urlcolor     = blue, 
  linkcolor    = blue, 
  citecolor    = red 
}

% Comments
\usepackage{comment}
\usepackage{todonotes}

% References
\usepackage{cite}

%%%%%%%%%%%%%%%%%%%%%%%%%%%%%%%%%%%%%%%%%%%%%%%%%%%%%%%%%%%%%%%%%%%%%%%%%%%%%%%%

\begin{document}

\title{Magnetoencephalography (MEG) Data Classification: Intra- vs. Cross-Subject Analysis}

\author{\IEEEauthorblockN{1\textsuperscript{st} Kian Niek Hamidi Abd Abad}
\IEEEauthorblockA{\textit{Utrecht University} \\
Utrecht, Netherlands \\
0609927}
\and
\IEEEauthorblockN{2\textsuperscript{nd} Márton Mohácsi}
\IEEEauthorblockA{\textit{Utrecht University} \\
Utrecht, Netherlands \\
1621106}
\and
\IEEEauthorblockN{3\textsuperscript{rd} Sid Pranjale}
\IEEEauthorblockA{\textit{Utrecht University} \\
Utrecht, Netherlands \\
4368959}
\and
\IEEEauthorblockN{4\textsuperscript{th} Jorik Roebroek}
\IEEEauthorblockA{\textit{Utrecht University} \\
Utrecht, Netherlands \\
9490466}
}

\maketitle

\begin{abstract}
% Briefly summarize the project goals, the MEG dataset, the deep learning architecture chosen, and the primary findings regarding intra-subject vs. cross-subject classification performance.
Your abstract text goes here.
\end{abstract}

\begin{IEEEkeywords}
Deep Learning, Magnetoencephalography (MEG), Brain Decoding, Time-Series Classification
\end{IEEEkeywords}

\section{Introduction}
% Introduce MEG data and the significance of inferring mental states from brain activity. State the goal: classifying whether a subject is in a resting, math, memory, or motor state.
Insert introduction text here.

\section{Data and Preprocessing}
% Detail the dataset structure (248 sensors x 35624 time steps).
% Discuss the preprocessing steps you took as suggested in the hints.

\subsection{Downsampling}
% Explain how and why you downsampled from the original 2034 Hz to improve training speed.

\subsection{Scaling and Normalization}
% Detail your min-max scaling or Z-score normalization strategy, especially focusing on time-wise scaling to adjust the 10e-15 fT magnitude.

\section{Methodology}

\subsection{Model Architecture (Task A)}
% Address Task (a) here: Choose a suitable deep learning model (e.g., CNN, LSTM, or a hybrid) for classifying the MEG time-series data. Justify your choice based on the data's temporal and spatial characteristics.
Detail your chosen deep learning model and justification here.

\subsection{Hyper-parameter Selection (Task C)}
% Address Task (c) here: Explain the choices of hyperparameters (learning rate, batch size, number of layers, etc.) and analyze their influence on the results for both classification types. How were they selected?
Detail your hyperparameter tuning approach.

\section{Experimental Setup}
% Explain how the data was split and how you handled memory management during training (e.g., batch loading the 64 files).

\subsection{Intra-Subject Classification}
% Describe the setup for training and testing on the same subject.

\subsection{Cross-Subject Classification}
% Describe the setup for training on a set of subjects and testing on unseen subjects.

\section{Results and Discussion}

\subsection{Performance Comparison (Task B)}
% Address Task (b) here: Compare the accuracy of intra-subject and cross-subject classification. Include visualizations such as bar charts or confusion matrices.
Explain your results and the differences in accuracy.

\begin{figure} [h]
    \centering
    %\includegraphics[width=0.8\linewidth]{results_graph.png}
    \caption{Comparison of training and validation accuracies across Intra-Subject and Cross-Subject models.}
    \label{fig:results_label}
\end{figure}

\subsection{Training vs. Testing Analysis (Task D)}
% Address Task (d) here: Discuss any significant differences in training and testing accuracies (e.g., signs of overfitting, especially in cross-subject). 
Analyze the performance gaps here.

\subsection{Alternative Approaches (Task D)}
% Address the second part of Task (d): Propose alternative models or techniques (e.g., regularization, different architectures) to improve results. Select one, implement it, and justify the choice.
Detail your selected alternative approach and its impact on the results.

\section{Conclusion}
% Summarize the findings of your MEG classification project, the effectiveness of the chosen models, and potential future work.
Summary of findings and future work.

\nocite{*}
\bibliographystyle{IEEEtran}
\bibliography{references}

\end{document}